\solchapter{em}

\sol{em}{11} Laplace's equation implies the mean-value property: $V$ at a point equals its average over any small sphere around it. Can a function with that property have a strict minimum?

\sol{em}{12} The positive charges move with the test charge, so they are at rest in its frame. Find the negative charges' speed in that frame with the velocity-addition formula, and compare how the two densities change, to order $v^2/c^2$.

\solchapter{thermo}

\sol{thermo}{10} How much can the demon lower the gas's entropy for each bit it learns? Compare that with the minimum cost of erasing each bit.

\sol{thermo}{11} For radiation $U = u(T)V$, so $(\partial U/\partial V)_T = u$, and $P = u/3$ gives $(\partial P/\partial T)_V = u'(T)/3$. The identity becomes a differential equation for $u(T)$.

\solchapter{optics}

\sol{optics}{10} A stationary deviation needs $\dd r/\dd i = \tfrac12$. Differentiate Snell's law to get $\dd r/\dd i$ in terms of $i$ and $r$, then eliminate $r$.

\sol{optics}{11} Divide the open wavefront beyond the disc's edge into Fresnel zones, rings whose distances to the screen point increase by half a wavelength each. Their contributions alternate in sign and shrink slowly. What is the sum?

\sol{optics}{12} Optics is reversible: light can retrace any path. If the target were hotter than the Sun, which way would net radiation flow along those paths?

\solchapter{statmech}

\sol{statmech}{11} Solve $n\lambda_T^3 = 2.6$ for $\lambda_T$, then $\lambda_T = h/\sqrt{2\pi m\kB T}$ for $T$. For liquid helium, get $n$ from the mass density.

\sol{statmech}{12} Use Stirling's approximation for $\ln\binom Nn$ and differentiate with respect to $n$. For the heat flow, ask which direction of energy transfer increases the total entropy when one of the two temperatures is negative.

\sol{statmech}{13} Compare $S(2N,2V,T)$ with $2S(N,V,T)$. Then subtract $\kB\ln N!$ using Stirling, and watch what happens to the $\ln V$.

\solchapter{sr}

\sol{sr}{9} The turnaround is the event $(5\un{yr}, 4\un{ly})$ in the Earth's frame. Find where the traveller's lines of simultaneity through that event --- one for the outgoing frame, one for the returning frame --- cross the Earth's worldline.

\sol{sr}{10} Are the two door-closings simultaneous in the pole's frame? For the shut door: how quickly can the news that the front has stopped reach the back?

\sol{sr}{11} Note that $c^2/g\approx1\un{ly}$ and $c/g\approx1\un{yr}$. Solve for the proper time to reach the midpoint and double it. For the catch, think about the rocket equation with exhaust at the speed of light.
